\thispagestyle{fancy}
\fancyhead[R]{A.U 2025-2026}
\fancyhead[L]{ chapitre 2: Analyse et spécification des besoins }
\chapter{Analyse et spécification des besoins}
\section*{Introduction}
Dans ce chapitre, nous analysons en détail les besoins de la plateforme Sellio. Nous identifions d'abord les acteurs, puis les besoins fonctionnels et non fonctionnels. Nous présentons ensuite les diagrammes de cas d'utilisation, l'architecture retenue, le backlog produit et la planification des sprints, et enfin l'environnement de travail.

\section{Analyse et spécification des besoins}
La spécification des besoins fonctionnels et non fonctionnels est la première étape du cycle de développement : elle fixe ce que la plateforme doit offrir et les qualités qu'elle doit garantir.

\subsection{Identification des acteurs}
Sellio met en relation une plateforme, des commerçants et les clients de ces commerçants. Chaque acteur y joue un rôle distinct.
\subsubsection*{\underline{\textbf{Acteurs principaux}}}
\begin{itemize}[label=\ding{109}]
\item \textbf{Administrateur Sellio :} membre de l'équipe de la plateforme. Il consulte les statistiques globales, la liste des boutiques et de leurs propriétaires, valide ou refuse les dossiers de vérification (KYC), suspend une boutique, bloque ou supprime un commerçant. Il n'a jamais accès aux clients, aux commandes ni aux analyses des boutiques.
\item \textbf{Super administrateur Sellio :} administrateur Sellio qui, en plus, ajoute des administrateurs, change le rôle des comptes de la plateforme et supprime un compte Sellio.
\item \textbf{Commerçant :} propriétaire d'une boutique. Il s'inscrit, crée sa boutique, soumet son dossier de vérification, puis gère depuis le back-office son catalogue, sa vitrine, ses commandes, ses retours, ses promotions, son programme de fidélité, ses campagnes e-mail et consulte l'analyse du comportement de ses visiteurs.
\item \textbf{Collaborateur de boutique :} membre de l'équipe d'un commerçant, doté d'un rôle aux permissions limitées (par exemple gestion des commandes uniquement).
\item \textbf{Visiteur :} internaute non connecté qui parcourt la vitrine d'une boutique, recherche des produits, reçoit des recommandations, remplit son panier et peut utiliser la cabine d'essayage virtuel.
\item \textbf{Client :} visiteur inscrit sur une boutique ; son compte n'existe que dans cette boutique, et la même personne peut avoir un autre compte, avec un autre mot de passe, dans une autre boutique. En plus des actions du visiteur, il passe commande, paie en ligne, suit ses commandes, demande un retour, publie des avis et cumule des points de fidélité.
\end{itemize}
\subsubsection*{\underline{\textbf{Acteurs secondaires}}}
\begin{itemize}[label=\ding{109}]
\item \textbf{Stripe :} prestataire de paiement. Il traite les paiements des commandes (\textit{PaymentIntent}) et l'enregistrement de la carte du commerçant lors de la vérification (\textit{SetupIntent}).
\item \textbf{Decart :} service d'intelligence artificielle générative qui produit en temps réel la vidéo d'essayage virtuel à partir du flux caméra du client et de l'image du vêtement.
\item \textbf{Brevo~\cite{brevo} :} service d'envoi d'e-mails transactionnels et marketing (codes de vérification, réinitialisation du mot de passe, invitations, confirmations de commande, newsletters).
\item \textbf{GitHub Actions et Dependabot :} GitHub Actions construit et teste la plateforme à chaque modification. Dependabot propose les mises à jour de dépendances. Aucun des deux ne déploie Sellio sur un serveur.
\item \textbf{Moteur d'apprentissage automatique :} ensemble de scripts Python (scikit-learn) appelés par le microservice d'analyse pour entraîner les modèles et produire recommandations, segments et scores d'attrition.
\end{itemize}

\subsection{Spécification des besoins}
Nous classons les besoins en deux familles : les besoins fonctionnels, qui décrivent \textbf{ce que} le système doit faire, et les besoins non fonctionnels, qui décrivent \textbf{comment} il doit le faire.
\subsubsection*{\underline{\textbf{Besoins fonctionnels}}}
Les besoins sont regroupés par module, ce qui reflète le découpage de la plateforme en microservices. Pour chaque fonctionnalité, nous précisons le ou les acteurs concernés.

\Needspace{20\baselineskip}
\paragraph{Module Plateforme, authentification et vérification}
Le tableau~\ref{tab:besoins-plateforme} présente les besoins de ce module.
\begin{table}[H]
\centering
\renewcommand{\arraystretch}{1.3}
\begin{tabular}{|p{9.5cm}|p{5.5cm}|}
\hline
\textbf{Fonctionnalité} & \textbf{Acteur(s) concerné(s)} \\
\hline
S'inscrire et s'authentifier de façon sécurisée (JWT et jeton de rafraîchissement) & Commerçant, Client, Administrateurs Sellio \\
\hline
Confirmer son adresse par un code envoyé par e-mail ; réinitialiser son mot de passe & Commerçant, Client, Brevo \\
\hline
Avoir un compte client distinct dans chaque boutique & Client \\
\hline
Créer une boutique (nom, secteur, identité visuelle) & Commerçant \\
\hline
Soumettre un dossier de vérification : carte bancaire, pièce d'identité (CIN ou passeport) et selfie & Commerçant \\
\hline
Valider ou refuser un dossier de vérification avec motif & Administrateur Sellio \\
\hline
Consulter les statistiques, les boutiques et leurs propriétaires & Administrateur Sellio \\
\hline
Suspendre une boutique, bloquer un commerçant & Administrateur Sellio \\
\hline
Supprimer un commerçant avec sa boutique et toutes ses données & Administrateur Sellio \\
\hline
Ajouter un administrateur, changer le rôle d'un compte de la plateforme & Super administrateur Sellio \\
\hline
Gérer les rôles et permissions de l'équipe de la boutique & Commerçant \\
\hline
\end{tabular}
\caption{Besoins fonctionnels du module Plateforme, authentification et vérification}
\label{tab:besoins-plateforme}
\end{table}

\Needspace{20\baselineskip}
\paragraph{Module Catalogue et vitrine}
Le tableau~\ref{tab:besoins-catalogue} présente les besoins de ce module.
\begin{table}[H]
\centering
\renewcommand{\arraystretch}{1.3}
\begin{tabular}{|p{9.5cm}|p{5.5cm}|}
\hline
\textbf{Fonctionnalité} & \textbf{Acteur(s) concerné(s)} \\
\hline
Gérer les produits, leurs images et leurs variantes & Commerçant, Collaborateur \\
\hline
Renseigner une fiche adaptée au secteur et au type du produit (bague, montre, chaussures, machine\dots) & Commerçant, Collaborateur \\
\hline
Choisir les valeurs dans des listes normalisées et les étendre & Commerçant, Collaborateur \\
\hline
Gérer les catégories & Commerçant, Collaborateur \\
\hline
Choisir le secteur et l'un des huit modèles de vitrine & Commerçant \\
\hline
Personnaliser l'apparence : logo et taille, couleurs, polices, barre d'annonce, coordonnées, réseaux sociaux & Commerçant \\
\hline
Personnaliser la page d'accueil (sections, ordre, textes, produits mis en avant) avec aperçu en direct & Commerçant \\
\hline
Parcourir la vitrine, filtrer et rechercher les produits & Visiteur, Client \\
\hline
\end{tabular}
\caption{Besoins fonctionnels du module Catalogue et vitrine}
\label{tab:besoins-catalogue}
\end{table}

\Needspace{20\baselineskip}
\paragraph{Module Transactions}
Le tableau~\ref{tab:besoins-transactions} présente les besoins de ce module.
\begin{table}[H]
\centering
\renewcommand{\arraystretch}{1.3}
\begin{tabular}{|p{9.5cm}|p{5.5cm}|}
\hline
\textbf{Fonctionnalité} & \textbf{Acteur(s) concerné(s)} \\
\hline
Gérer un panier (taille obligatoire pour les articles concernés) & Visiteur, Client \\
\hline
Passer commande avec livraison, TVA et code promo & Client \\
\hline
Payer en ligne par carte  & Client, Stripe \\
\hline
Suivre ses commandes et demander un retour & Client \\
\hline
Gérer les commandes, leurs statuts et les retours & Commerçant, Collaborateur \\
\hline
Paramétrer la TVA, les zones de livraison et la politique de retour & Commerçant \\
\hline
Publier un avis ; modérer et répondre aux avis & Client ,Commerçant \\
\hline
\end{tabular}
\caption{Besoins fonctionnels du module Transactions}
\label{tab:besoins-transactions}
\end{table}

\Needspace{20\baselineskip}
\paragraph{Module Marketing et fidélisation}
Le tableau~\ref{tab:besoins-marketing} présente les besoins de ce module.
\begin{table}[H]
\centering
\renewcommand{\arraystretch}{1.3}
\begin{tabular}{|p{9.5cm}|p{5.5cm}|}
\hline
\textbf{Fonctionnalité} & \textbf{Acteur(s) concerné(s)} \\
\hline
Créer des codes promo et des remises automatiques & Commerçant \\
\hline
Publier des bannières ciblées par segment & Commerçant \\
\hline
Envoyer une newsletter ou contacter un client par e-mail & Commerçant, Brevo \\
\hline
Configurer le programme de fidélité et ses niveaux & Commerçant \\
\hline
Consulter son solde de points et son niveau & Client \\
\hline
\end{tabular}
\caption{Besoins fonctionnels du module Marketing et fidélisation}
\label{tab:besoins-marketing}
\end{table}

\Needspace{20\baselineskip}
\paragraph{Module Intelligence (essayage virtuel et apprentissage automatique)}
Le tableau~\ref{tab:besoins-ia} présente les besoins de ce module.
\begin{table}[H]
\centering
\renewcommand{\arraystretch}{1.3}
\begin{tabular}{|p{9.5cm}|p{5.5cm}|}
\hline
\textbf{Fonctionnalité} & \textbf{Acteur(s) concerné(s)} \\
\hline
Essayer virtuellement un vêtement en temps réel avec sa caméra & Visiteur, Client, Decart \\
\hline
Collecter les événements de navigation & Visiteur, Client (implicite) \\
\hline
Recevoir des recommandations : produits similaires, pour vous, souvent achetés ensemble, populaires & Visiteur, Client, Moteur ML \\
\hline
Consulter l'analyse comportementale : entonnoir, segments de visiteurs, produits consultés & Commerçant \\
\hline
Identifier les clients à risque d'attrition & Commerçant, Moteur ML \\
\hline
Ré-entraîner les modèles de la boutique & Commerçant, Moteur ML \\
\hline
\end{tabular}
\caption{Besoins fonctionnels du module Intelligence}
\label{tab:besoins-ia}
\end{table}

\Needspace{16\baselineskip}
\paragraph{Module Qualité et déploiement}
Le tableau~\ref{tab:besoins-devops} présente les besoins de ce module.
\begin{table}[H]
\centering
\renewcommand{\arraystretch}{1.3}
\begin{tabular}{|p{9.5cm}|p{5.5cm}|}
\hline
\textbf{Fonctionnalité} & \textbf{Acteur(s) concerné(s)} \\
\hline
Construire et tester automatiquement chaque modification & Équipe de développement , GitHub Actions \\
\hline
Construire les images Docker une fois les tests passés, sans les déployer & GitHub Actions \\
\hline
Lancer toute la plateforme en une commande & Équipe de développement  \\
\hline
Proposer les mises à jour de dépendances & Dependabot \\
\hline
\end{tabular}
\caption{Besoins fonctionnels du module Qualité et déploiement}
\label{tab:besoins-devops}
\end{table}

\subsubsection*{\underline{\textbf{Besoins non fonctionnels}}}
Les besoins non fonctionnels définissent les qualités attendues du système. \newline
\Needspace{20\baselineskip}
Le tableau~\ref{tab:exigences-non-fonctionnelles} les présente selon le modèle de qualité ISO/IEC 25010~\cite{ISO25010}.

\begin{longtable}{|p{3.2cm}|p{4cm}|p{8cm}|}
\caption{Exigences non fonctionnelles du système}
\label{tab:exigences-non-fonctionnelles} \\
\hline
\textbf{Caractéristique} & \textbf{Sous-caractéristique} & \textbf{Description} \\
\hline
\endfirsthead
\hline
\textbf{Caractéristique} & \textbf{Sous-caractéristique} & \textbf{Description} \\
\hline
\endhead
\hline \multicolumn{3}{r}{\textit{Suite à la page suivante}} \\
\endfoot
\hline
\endlastfoot

\multirow{5}{*}{Sécurité} & Authentification & Jeton d'accès JWT (24 h) et jeton de rafraîchissement (7 jours) ; mots de passe hachés avec BCrypt. \\
\cline{2-3}
 & Contrôle d'accès & Rôles et permissions vérifiés côté serveur ; un commerçant n'accède qu'aux données de sa propre boutique ; l'équipe Sellio n'accède à aucune donnée de client. \\
\cline{2-3}
 & Protection des paiements & Aucune donnée de carte ne transite par nos serveurs ; seuls la marque, les 4 derniers chiffres, l'expiration et la référence Stripe sont stockés (PCI-DSS). \\
\cline{2-3}
 & Confiance & Une boutique reste « en attente » tant que son dossier KYC n'est pas validé ; un compte bloqué ne peut plus se connecter. \\
\cline{2-3}
 & Secrets & Aucune clé dans le code ni dans les images : elles sont lues au démarrage depuis un fichier \texttt{.env} exclu du dépôt. \\
\hline
\multirow{2}{*}{Performance} & Temps de réponse & Recommandations servies depuis des modèles pré-entraînés ; les événements sont envoyés par lots pour ne pas ralentir la navigation. \\
\cline{2-3}
 & Volume de données & Table d'événements partitionnée par mois, index BRIN sur la date, agrégats quotidiens et rétention de 13 mois. \\
\hline
\multirow{2}{*}{Fiabilité} & Tolérance aux pannes & La panne du service d'analyse n'empêche pas les achats : la vitrine masque les recommandations. En revanche, la réservation de stock est un appel synchrone au catalogue : s'il est arrêté, la commande ne peut pas être enregistrée. \\
\cline{2-3}
 & Intégrité & Transactions ACID pour les commandes et les paiements. \\
\hline
\multirow{2}{*}{Utilisabilité} & Ergonomie & Interfaces responsives ; thème sombre par défaut sur l'espace Sellio avec bascule vers le thème clair. \\
\cline{2-3}
 & Personnalisation & Toute modification du thème est visible immédiatement sur la vitrine, sans écrire de code. \\
\hline
\multirow{3}{*}{Maintenabilité} & Modularité & Un microservice par domaine métier, une bibliothèque partagée pour les entités et DTO. \\
\cline{2-3}
 & Évolutivité & Ajout d'un modèle de vitrine, d'un secteur, d'un type de produit ou d'un modèle ML sans toucher aux autres services. \\
\cline{2-3}
 & Testabilité & Tests automatisés (JUnit, Vitest, pytest) exécutés à chaque modification. \\
\hline
\multirow{2}{*}{Portabilité} & Compatibilité & Navigateurs modernes (Chrome, Firefox, Edge) ; l'essayage virtuel exige une caméra et WebRTC. \\
\cline{2-3}
 & Déploiement & Chaque service est livré sous forme d'image Docker ; la plateforme se lance en une commande (Docker Compose). \\
\hline
Confidentialité & Données personnelles & Le suivi comportemental repose sur un identifiant de visiteur anonyme ; les pièces d'identité ne sont visibles que par l'administrateur Sellio. \\
\hline
\end{longtable}

\subsection{Diagramme de cas d'utilisation général}
Le système est utilisé par six acteurs principaux : \textbf{le Super administrateur} et \textbf{l'Administrateur Sellio}, \textbf{le Commerçant}, \textbf{le Collaborateur de boutique}, \textbf{le Visiteur} et \textbf{le Client}. Les sous-sections suivantes présentent les cas d'utilisation de chacun.

\subsubsection{Cas d'utilisation de l'équipe Sellio}
La figure~\ref{fig:cu-admin} présente les cas d'utilisation des administrateurs de la plateforme ; le super administrateur hérite de tous les cas de l'administrateur.
\begin{figure}[H]
    \centering
    \img{0.9}{cu_admin_sellio.png}
    \caption{Diagramme de cas d'utilisation de l'équipe Sellio}
    \label{fig:cu-admin}
\end{figure}

\subsubsection{Cas d'utilisation du Commerçant}
La figure~\ref{fig:cu-commercant} présente les cas d'utilisation du commerçant. Le collaborateur de boutique hérite d'une partie de ces cas, selon les permissions de son rôle.
\begin{figure}[H]
    \centering
    \img{0.7}{cu_commercant.png}
    \caption{Diagramme de cas d'utilisation du Commerçant}
    \label{fig:cu-commercant}
\end{figure}

\subsubsection{Cas d'utilisation du Visiteur et du Client}
La figure~\ref{fig:cu-client} présente les cas d'utilisation du visiteur et du client. Le client hérite de tous les cas du visiteur.
\begin{figure}[H]
    \centering
    \img{0.9}{cu_client.png}
    \caption{Diagramme de cas d'utilisation du Visiteur et du Client}
    \label{fig:cu-client}
\end{figure}

\section{Analyse globale}
\subsection{Diagramme de classes d'analyse}
La figure~\ref{fig:classes} présente le diagramme de classes d'analyse. La classe centrale est \texttt{Shop} : chaque utilisateur, produit, catégorie, commande, promotion ou événement comportemental est rattaché à une boutique, ce qui matérialise l'isolation multi-boutiques. Une boutique possède un dossier de vérification (\texttt{MerchantVerification}), un statut (\texttt{PENDING}, \texttt{ACTIVE}, \texttt{REJECTED}, \texttt{SUSPENDED}), un secteur, un modèle de vitrine, des couleurs (\texttt{theme}) et sa personnalisation (\texttt{settings}, objet JSON : logo, coordonnées, page d'accueil\dots). Le diagramme présente le modèle du domaine dans son ensemble ; physiquement, ces classes sont réparties entre les bases des microservices (tableau~\ref{tab:bases}), et les associations qui franchissent une frontière de service (une commande et son client, un avis et son auteur) sont matérialisées par un identifiant. Les caractéristiques propres au secteur d'un produit sont rangées dans \texttt{Product.attributes}, ce qui évite une table par secteur. Les codes envoyés par e-mail (\texttt{AuthCode}) ne sont stockés que sous forme d'empreinte.
\begin{figure}[H]
    \centering
    \img{1}{diagramme_classes.png}
    \caption{Diagramme de classes d'analyse}
    \label{fig:classes}
\end{figure}

\subsection{Architecture globale}
L'architecture influence directement les performances, la maintenabilité et l'évolutivité de la plateforme. Trois approches ont été comparées : \textbf{l'architecture monolithique}, \textbf{l'architecture N-tiers} et \textbf{l'architecture microservices}.\newline
\Needspace{20\baselineskip}
Le tableau~\ref{tab:comparaison_architectures} présente leurs avantages et leurs limites.

\begin{longtable}{|>{\raggedright\arraybackslash}p{3.2cm}|>{\raggedright\arraybackslash}p{11cm}|}
\hline
\textbf{Architecture} & \textbf{Caractéristiques} \\
\hline
\endfirsthead
\hline
\textbf{Architecture} & \textbf{Caractéristiques} \\
\hline
\endhead
Monolithique &
\begin{itemize}[leftmargin=*, nosep]
    \item Présentation, logique métier et accès aux données dans une seule unité de déploiement.
    \item \textbf{Avantages} : simplicité de développement et de déploiement, débogage facile, cohérence transactionnelle native.
    \item \textbf{Limites} : mise à l'échelle uniquement globale, un défaut dans un module peut affecter tout le système, maintenance difficile quand le code grossit.
\end{itemize} \\
\hline
Client--serveur N-tiers &
\begin{itemize}[leftmargin=*, nosep]
    \item Séparation en couches : présentation, logique métier, données.
    \item \textbf{Avantages} : séparation claire des responsabilités, plusieurs clients peuvent consommer la même couche métier.
    \item \textbf{Limites} : le serveur d'application reste une unité de déploiement unique, avec un couplage interne entre modules.
\end{itemize} \\
\hline
Microservices &
\begin{itemize}[leftmargin=*, nosep]
    \item Décomposition en services indépendants, chacun responsable d'un domaine métier, déployés séparément et communiquant par API REST derrière une passerelle.
    \item \textbf{Avantages} : mise à l'échelle service par service, isolation des pannes, déploiements indépendants, liberté technologique par service.
    \item \textbf{Limites} : complexité opérationnelle, latence réseau entre services, cohérence des données distribuées plus difficile.
\end{itemize} \\
\hline
\caption{Comparaison des architectures logicielles étudiées}
\label{tab:comparaison_architectures}
\end{longtable}

L'\textbf{architecture microservices} a été retenue, pour trois raisons. D'abord, le sujet lui-même demande une plateforme distribuée. Ensuite, les domaines de Sellio ont des charges très différentes : le service d'analyse reçoit un flux continu d'événements et lance des entraînements de modèles coûteux, alors que le service d'identité traite peu de requêtes ; les séparer permet de dimensionner chacun selon sa charge et d'éviter qu'un entraînement ne ralentisse les achats. Enfin, le projet est né d'un monolithe (NaturEssence) dont la croissance rendait les évolutions de plus en plus risquées ; la migration vers des services séparés a rendu chaque domaine plus facile à faire évoluer.


\Needspace{18\baselineskip}
\subsubsection{Une base de données par microservice}
Le découpage n'est pas seulement logique. Chaque microservice possède \textbf{son propre serveur PostgreSQL}, seul processus autorisé à ouvrir cette base (\textit{database per service}). Il ne lit jamais les tables d'un autre : il demande l'information par une API interne. Sauvegarder, migrer ou saturer un disque ne concerne qu'un domaine. Le tableau~\ref{tab:bases} précise le contenu de chaque serveur.
\begin{table}[H]
\centering
\renewcommand{\arraystretch}{1.3}
\small
\begin{tabular}{|p{3.1cm}|p{2.3cm}|p{9.2cm}|}
\hline
\textbf{Base} & \textbf{Propriétaire} & \textbf{Données} \\
\hline
\texttt{sellio\_auth} & auth-service & Comptes, rôles et permissions, segments, boutiques et leur personnalisation, dossiers KYC, codes e-mail, jetons de rafraîchissement, fidélité (configuration et points) \\
\hline
\texttt{sellio\_catalog} & catalog-service & Produits, variantes, catégories, collections, avis \\
\hline
\texttt{sellio\_orders} & order-service & Commandes et lignes, retours et politique de retour, codes promo et utilisations, remises, zones de livraison, TVA \\
\hline
\texttt{sellio\_marketing} & marketing-service & Bannières, apparence par défaut de Sellio \\
\hline
\texttt{sellio\_analytics} & analytics-service & Événements (partitionnés par mois), agrégats quotidiens, recommandations, segments, historique des entraînements \\
\hline
\end{tabular}
\caption{Répartition des données entre les bases des microservices}
\label{tab:bases}
\end{table}

Ce découpage a trois conséquences sur la conception :
\begin{itemize}
  \item \textbf{Références par identifiant.} Une commande ne pointe pas vers l'entité client d'un autre service : elle conserve l'identifiant du client et une copie des informations utiles à son affichage (nom, e-mail, adresse ; nom et prix de chaque article au moment de l'achat).
  \item \textbf{API internes entre services.} Les échanges passent par des routes \texttt{/internal/**} protégées par un jeton de service et non exposées par la passerelle : le service de commandes demande au catalogue de réserver le stock et au service d'identité de créditer les points de fidélité ; le service d'analyse demande chaque nuit au catalogue les produits et au service de commandes les commandes payées.
  \item \textbf{Boutique connue de tous sans lecture de la base d'identité.} Le jeton JWT porte l'identifiant et le \textit{slug} de la boutique de l'utilisateur ; pour les requêtes publiques de la vitrine, la correspondance \textit{slug} / boutique est obtenue auprès du service d'identité et mise en cache.
\end{itemize}

Ces cinq serveurs tournent aujourd'hui sur le même hôte, pour le développement et la démonstration, mais ce sont cinq processus, cinq volumes et cinq chaînes de connexion. Les déplacer sur cinq machines ne change que ces adresses. La bibliothèque \texttt{shared-lib} ne porte que les contrats échangés (DTO et vérification du jeton) : le schéma de chaque base reste dans son service.

\subsubsection{Choix du système de gestion de base de données}
 Le tableau~\ref{tab:choix-sgbd} résume la comparaison qui nous a conduits à retenir PostgreSQL.
\begin{table}[H]
\centering
\renewcommand{\arraystretch}{1.4}
\begin{tabular}{|p{4cm}|p{5.3cm}|p{5.3cm}|}
\hline
\textbf{Critère} & \textbf{MongoDB} & \textbf{PostgreSQL} \\
\hline
Transactions (commande + paiement + stock) & Multi-documents possibles mais plus coûteuses & ACID natives \\
\hline
Intégrité référentielle (boutique, client, commande) & À gérer dans l'application & Clés étrangères et contraintes \\
\hline
Données flexibles (attributs, thème) & Natif & Type JSONB \\
\hline
Volume d'événements & Bon & Partitionnement déclaratif et index BRIN \\
\hline
Requêtes analytiques (agrégats, jointures) & Pipeline d'agrégation & SQL standard, fonctions de fenêtre \\
\hline
Une base par service sur un même serveur & Oui & Oui (bases isolées, rôles distincts) \\
\hline
\end{tabular}
\caption{Comparaison MongoDB / PostgreSQL}
\label{tab:choix-sgbd}
\end{table}
Les données de Sellio sont avant tout relationnelles et transactionnelles (une commande relie un client, des produits, une boutique et un paiement) ; PostgreSQL couvre en outre le besoin de flexibilité (JSONB) et de volume (partitionnement), d'où ce choix~\cite{postgresql2026}.

\subsubsection{Architecture logique de la plateforme}
La figure~\ref{fig:archi_logique} présente l'architecture logique de Sellio et ses couches.
\begin{figure}[H]
\centering
\img{1}{architecture_logique.png}
\caption{Architecture logique de la plateforme Sellio}
\label{fig:archi_logique}
\end{figure}

\subsubsection*{La couche présentation (front-end) :}
Elle se compose de deux applications React :
\begin{itemize}
  \item \textbf{l'application Sellio} (port 3000), qui regroupe le site de la plateforme (accueil, inscription, connexion), la console de l'administrateur Sellio et le back-office de chaque commerçant ;
  \item \textbf{la vitrine} (port 3001), commune à toutes les boutiques : la boutique affichée est déterminée par son identifiant (\textit{slug}) ; son modèle, ses couleurs, ses polices et sa page d'accueil sont chargés dynamiquement. Les éditeurs du back-office l'affichent dans un cadre pour en montrer l'aperçu en direct.
\end{itemize}

\subsubsection*{La couche passerelle :}
Une passerelle \textbf{Spring Cloud Gateway} (port 8080) est le point d'entrée unique des API. Elle route chaque requête vers le bon microservice selon son chemin (\texttt{/api/v1/auth/**} vers le service d'identité, \texttt{/api/v1/public/products/**} vers le catalogue, etc.).

\subsubsection*{La couche métier (back-end) :}
Elle est composée de cinq microservices Spring Boot :
\begin{itemize}
  \item \textbf{auth-service} (8081) : comptes, JWT, rôles et permissions, boutiques et leur personnalisation, vérification KYC, console plateforme, fidélité et segments ;
  \item \textbf{catalog-service} (8082) : produits et leurs caractéristiques par secteur, variantes, catégories, avis, téléversement d'images et jeton d'essayage virtuel ;
  \item \textbf{order-service} (8083) : tunnel de commande, paiement Stripe, commandes, retours, TVA, livraison, promotions et tableau de bord ;
  \item \textbf{marketing-service} (8084) : bannières et e-mailing ;
  \item \textbf{analytics-service} (8085) : collecte des événements, recommandations, analyse comportementale et prédiction de l'attrition.
\end{itemize}

\subsubsection*{La couche intelligence :}
Des scripts \textbf{Python} (pandas, scikit-learn) entraînent et appliquent les modèles. Le service d'analyse les exécute comme sous-processus et stocke les modèles entraînés par boutique.

\subsubsection*{La couche données :}
Chaque microservice a \textbf{son propre serveur PostgreSQL}, dans un conteneur distinct, avec son port, son compte et son volume : \texttt{postgres-auth} (5432), \texttt{postgres-catalog} (5433), \texttt{postgres-orders} (5434), \texttt{postgres-marketing} (5435) et \texttt{postgres-analytics} (5436). Un service ne connaît pas l'adresse des autres bases. Les images de produits sont sur un volume du catalogue.

\subsubsection*{Les services externes :}
\textbf{Stripe} (paiements), \textbf{Decart} (essayage virtuel) et \textbf{Brevo} (e-mails).

\subsubsection{Architecture physique de la plateforme}
La figure~\ref{fig:archi_physique} présente le déploiement physique, sur un même hôte Docker Compose, mais avec un conteneur par processus. Il y a les deux front-ends nginx, la passerelle, les cinq microservices et \textbf{cinq serveurs PostgreSQL}. Arrêter ou remplir le disque de \texttt{postgres-analytics} ne rend pas les comptes ni les commandes illisibles : chaque volume est indépendant. Les services ne se parlent que par API. Le navigateur parle en HTTPS aux front-ends et à la passerelle, et directement à Stripe (saisie de carte) et à Decart (WebRTC).
\begin{figure}[H]
    \centering
    \img{1}{architecture_physique.png}
    \caption{Architecture physique de la plateforme Sellio}
    \label{fig:archi_physique}
\end{figure}

\section{Pilotage du projet avec Scrum}
\subsection{Backlog du produit}

Le backlog produit regroupe les principales fonctionnalités de Sellio. 
Afin de conserver une vue globale du projet, les besoins sont ici synthétisés 
par grands modules. Les user stories détaillées sont présentées dans les sprints correspondants.

\begin{table}[H]
\centering
\renewcommand{\arraystretch}{1.3}
\begin{tabular}{|p{2.9cm}|p{9.5cm}|p{2.5cm}|}
\hline
\textbf{Module} & \textbf{Fonctionnalités principales} & \textbf{Complexité} \\
\hline

Authentification et utilisateurs &
Inscription, connexion, gestion des profils, rôles, permissions et réinitialisation du mot de passe. &
Élevée \\
\hline

Plateforme et boutiques &
Création de boutiques, gestion multi-boutiques, vérification KYC et console d'administration Sellio. &
Élevée \\
\hline

Catalogue et vitrine &
Gestion des produits et catégories, fiches adaptées aux huit secteurs, modèles de vitrine et personnalisation. &
Élevée \\
\hline

Parcours transactionnel &
Panier, commande, paiement Stripe, gestion des commandes, retours, TVA, livraison et avis. &
Élevée \\
\hline

Marketing et fidélisation &
Promotions, bannières, e-mailing, programme de fidélité et gestion des niveaux clients. &
Moyenne \\
\hline

Essayage virtuel &
Essayage de vêtements en temps réel à partir de la caméra grâce à l'API Decart. &
Élevée \\
\hline

Intelligence comportementale &
Collecte des événements, recommandations personnalisées, segmentation des visiteurs et prédiction de l'attrition. &
Élevée \\
\hline

Qualité et DevOps &
Tests automatisés, conteneurisation avec Docker et intégration continue avec GitHub Actions. &
Élevée \\
\hline

\end{tabular}
\caption{Synthèse du backlog produit de Sellio}
\label{tab:backlog-produit}
\end{table}

\subsection{Planification des releases et des sprints}
Le stage commence le 2 février 2026 et, par avenant, se termine le 29 octobre 2026. Les dix derniers jours, du 20 au 29 octobre, sont réservés aux tests automatisés, aux corrections du pipeline et à la mise au propre du rapport. Le développement tient donc du 2 février au 19 octobre, en douze sprints d'environ trois semaines. Cette durée, un peu plus longue qu'un sprint Scrum classique, correspond au rythme des revues avec l'encadrant et au fait qu'une seule personne réalise l'analyse, le code et les diagrammes.

\begin{table}[H]
\centering
\renewcommand{\arraystretch}{1.25}
\small
\begin{tabular}{|c|p{3.4cm}|c|p{7.2cm}|}
\hline
\textbf{Sprint} & \textbf{Dates} & \textbf{Release} & \textbf{Livrable} \\
\hline
1 & 2 févr. -- 22 févr. & 1 & Passerelle, authentification, rôles \\
\hline
2 & 23 févr. -- 15 mars & 1 & Boutiques, KYC, console Sellio \\
\hline
3 & 16 mars -- 5 avr. & 2 & Catalogue et fiches par secteur \\
\hline
4 & 6 avr. -- 26 avr. & 2 & Huit vitrines et personnalisation \\
\hline
5 & 27 avr. -- 17 mai & 3 & Panier, commande, paiement Stripe \\
\hline
6 & 18 mai -- 7 juin & 3 & Retours, TVA, livraison, avis \\
\hline
7 & 8 juin -- 28 juin & 4 & Promotions, bannières, e-mailing \\
\hline
8 & 29 juin -- 19 juil. & 4 & Fidélité et segments \\
\hline
9 & 20 juil. -- 16 août & 5 & Essayage virtuel Decart \\
\hline
10 & 17 août -- 13 sept. & 6 & Suivi comportemental et recommandation \\
\hline
11 & 14 sept. -- 4 oct. & 6 & Segmentation et attrition \\
\hline
12 & 5 oct. -- 19 oct. & 7 & Docker et pipeline d'intégration \\
\hline
Clôture & 20 oct. -- 29 oct. & --- & Tests, corrections de la CI, rapport \\
\hline
\end{tabular}
\caption{Calendrier des douze sprints, du 2 février au 29 octobre 2026}
\label{tab:calendrier-sprints}
\end{table}

Le sprint 9 est un peu plus long : l'essayage dépend d'un service externe (caméra, jeton, WebRTC) et les essais prennent plus de temps qu'un écran de gestion. Les sprints 10 et 11 aussi, parce que l'évaluation des modèles sur les jeux publics ne se réduit pas à l'écriture de l'API.

\section{Environnement de travail}
\subsection{Environnement matériel de développement}
Le développement a été réalisé sur un ordinateur portable présentant les caractéristiques suivantes :
\begin{itemize}
    \item \textbf{Modèle :} MSI GF63 Thin 10SC
    \item \textbf{Système d'exploitation :} Microsoft Windows 11 Professionnel (64 bits)
    \item \textbf{Processeur :} Intel Core i5-10500H @ 2.50 GHz, 6 cœurs et 12 processeurs logiques
    \item \textbf{Mémoire RAM :} 16 Go
   
\end{itemize}

\subsection{Environnement logiciel de développement}

Le tableau~\ref{tab:outils} présente les principaux outils utilisés.

\vspace{0.3em}

\renewcommand{\arraystretch}{1.3}
\begin{center}
\begin{tabular}{|l|p{9cm}|c|}
\hline
\textbf{Outil} & \textbf{Description} & \textbf{Logo} \\
\hline
Visual Studio Code~\cite{vscode} &
Éditeur principal pour le front-end React, le back-end Spring Boot et les scripts Python. &
\logo{25pt}{vscode.jpeg} \\
\hline
Git et GitHub~\cite{github} & Gestion de versions et hébergement du dépôt du projet. & \logo{25pt}{github.png} \\
\hline
PlantUML~\cite{plantuml} & Génération des diagrammes UML à partir de descriptions textuelles. & \logo{25pt}{plantuml.png} \\
\hline
PostgreSQL~\cite{postgresql2026} & Système de gestion de base de données relationnelle. & \logo{25pt}{postgresql.png} \\
\hline
Maven~\cite{maven} & Construction et gestion des dépendances des microservices Java. & \logo{25pt}{maven.png} \\
\hline
Stripe Dashboard~\cite{stripe} & Suivi des paiements et des cartes en mode test. & \logo{25pt}{stripe.png} \\
\hline
Docker~\cite{docker} & Images des services et des front-ends ; Docker Compose pour lancer la plateforme. & \logo{25pt}{docker.png} \\
\hline
GitHub Actions~\cite{githubactions} & Intégration continue : construction et tests à chaque modification. & \logo{25pt}{githubactions.png} \\
\hline
\end{tabular}
\captionof{table}{Outils utilisés pour la modélisation, le développement et le suivi}
\label{tab:outils}
\end{center}

\subsubsection{Langages de programmation}
\Needspace{16\baselineskip}
Le tableau~\ref{tab:langages} présente les langages utilisés.
\vspace{0.5em}
\begin{center}
\begin{tabular}{|l|p{9cm}|c|}
\hline
\textbf{Langage} & \textbf{Usage} & \textbf{Logo} \\
\hline
Java 17~\cite{java17} & Microservices (API REST, logique métier, sécurité). & \logo{25pt}{java.png} \\
\hline
JavaScript (JSX)~\cite{javascript} & Applications React : Sellio, back-office et vitrine. & \logo{25pt}{javascript.png} \\
\hline
Python 3~\cite{python} & Nettoyage des données, entraînement et évaluation des modèles. & \logo{25pt}{python.png} \\
\hline
SQL & Requêtes, partitionnement et agrégats PostgreSQL. & \logo{25pt}{sql.png} \\
\hline
\end{tabular}
\captionof{table}{Langages de programmation utilisés}
\label{tab:langages}
\end{center}

\subsubsection{Frameworks et bibliothèques}
\Needspace{24\baselineskip}
Le tableau~\ref{tab:frameworks} présente les frameworks et bibliothèques utilisés.
\vspace{0.5em}
\begin{center}
\begin{tabular}{|l|p{6cm}|c|}
\hline
\textbf{Framework} & \textbf{Rôle} & \textbf{Logo} \\
\hline
Spring Boot 3.2~\cite{springboot} & Socle des microservices (Spring Web, Spring Data JPA, Spring Security). & \logo{25pt}{springboot.png} \\
\hline
Spring Cloud Gateway~\cite{springcloudgateway} & Passerelle API : point d'entrée unique et routage vers les services. & \logo{25pt}{springcloud.png} \\
\hline
React 18~\cite{react} et Vite~\cite{vite} & Interfaces web à base de composants ; serveur de développement et compilation. & \logo{25pt}{react.png} \\
\hline
Tailwind CSS~\cite{tailwind} & Mise en forme utilitaire ; variables CSS pour les thèmes des vitrines. & \logo{25pt}{tailwind.png} \\
\hline
scikit-learn~\cite{scikitlearn}, pandas, NumPy & Modèles d'apprentissage automatique et préparation des données. & \logo{25pt}{scikitlearn.png} \\
\hline
Stripe.js~\cite{stripe} & Saisie sécurisée de la carte et confirmation des paiements. & \logo{25pt}{stripe.png} \\
\hline
SDK Decart~\cite{decart} & Connexion WebRTC au modèle d'essayage virtuel en temps réel. & \logo{25pt}{decart.png} \\
\hline
JUnit 5~\cite{junit}, Vitest~\cite{vitest}, pytest~\cite{pytest} & Tests automatisés du back-end, du front-end et des modèles. & \logo{25pt}{junit.png} \\
\hline
\end{tabular}
\captionof{table}{Frameworks et bibliothèques utilisés}
\label{tab:frameworks}
\end{center}

\section*{Conclusion}
Dans ce chapitre, nous avons identifié les acteurs, spécifié les besoins fonctionnels et non fonctionnels, justifié l'architecture microservices et le choix de PostgreSQL, puis présenté le backlog produit et la planification en sept releases et douze sprints. Le chapitre suivant présente la première release.
